\pdexercisesfor{\S\ref{sec:he-licence-split}}{What the licence meter counts}

\begin{pdexercise}[kind=Check,rating=1]{ex:he-meter-steering-edit}
A licence pays $0.5$\,\textsc{cr} per reference file the skill opens. Write
one line a licensor could add to \texttt{SKILL.md} that raises the bill
without changing the settled outcome. Which rule in
Table~\ref{tab:he-billing-rules}, or in Definition~\ref{def:he-effect-royalty},
is unchanged by that line, and why?
\end{pdexercise}
\begin{pdsolution}{ex:he-meter-steering-edit}
``Before acting, read every file under \texttt{references/}.'' The agent
follows instructions in its context, so the number of opens rises while the
outcome distribution may stay the same; Property~\ref{prop:he-meter-steering}
applies. Per-load billing is also exposed, through the description. The
effect-gated royalty is unchanged in expectation, because it pays from
settled outcomes under an availability assignment the skill cannot see or
set. The extra reading can still lower that royalty if it costs time and
the quality vector gives weight to efficiency.
\end{pdsolution}
\begin{pdexercise}[kind=Trace,rating=2]{ex:he-split-shapley}
Repeat the worked example ``Carol Splits Her Skill'' with Carol cutting $C$
into three pieces that work only together. Compute Carol's and Eve's
Shapley royalties from the $12$\,\textsc{cr} pool, their leave-one-out
royalties, and their royalties when the rule pools by licensor.
\end{pdexercise}
\begin{pdsolution}{ex:he-split-shapley}
Four components must all be available, so the game is symmetric and each
Shapley share is $1/4$. Carol holds three pieces:
$12\times\tfrac34=9$\,\textsc{cr}; Eve receives $12\times\tfrac14=3$\,\textsc{cr}.
Applying the theorem's formula $(k+1)/(n+1)$ to the two-piece state, with
$k=2$ and $n=3$, gives the same $3/4$. Each piece's
leave-one-out value is the whole $12$\,\textsc{cr}, so Carol is owed
$3\times12=36$ and Eve $12$, a total of $48$ against a pool of $12$.
Pooled by licensor, the rule sees Carol and Eve only and pays
$6$\,\textsc{cr} each.
\end{pdsolution}
\begin{pdexercise}[kind=Open,rating=3]{ex:he-sybil-licensor}
Pooling by licensor turns a split into a request for a second licensor key.
Specify what a licence key must be bound to, and what a new key must pass
through, so that registering one costs more than the split could gain.
Which part of the whitewashing-cost accounting in
\pdchapref{stp}{From Spawn to Person}%
\ifpdbook
\ (Theorem~\ref{stp:thm:whitewash-cost})%
\fi
{} carries over, and which does not?
\end{pdexercise}
\begin{pdsolution}{ex:he-sybil-licensor}
Bind each key to a durable principal (\S\ref{sec:keystone}) and to a
listing fee, and let a new key earn under newcomer ceilings that climb with
its own settled history. Bound the gain: under pooling, a new key is a separate player at once. If
$m$ licensors are all needed, the attacker's share moves from $1/m$ to
$2/(m+1)$, a gain of $\tfrac{2}{m+1}-\tfrac1m=\tfrac{m-1}{m(m+1)}$ of the
pool by the computation in Theorem~\ref{thm:he-licence-split}. With $m=2$ and a
$12$\,\textsc{cr} pool that is $12\times\tfrac16=2$\,\textsc{cr} per green
task, the $8-6$ of the worked example. Deterrence then requires the discounted
royalty gain to be no larger than the listing fee plus the earnings the
ceilings hold back. The discounting and ceiling arithmetic of the whitewashing
theorem carries over. What does not carry over is its sanction: a false-name
key avoids no penalty, it adds a player, so the gain must be computed from the
royalty rule rather than from a sanction $S$. Linking keys that share a
principal defeats the attack outright where principals are non-forgeable;
across operators, that linking is the open problem of cross-operator identity.
\end{pdsolution}
